% Generated draft: requires mathematical and source review.
\documentclass[11pt]{article}
\usepackage[margin=25mm]{geometry}
\usepackage{fontspec}
\setmainfont{Latin Modern Roman}
\usepackage{amsmath,amssymb,mathtools,unicode-math}
\setmathfont{Latin Modern Math}
\usepackage{enumitem,needspace,xurl,hyperref}
\setlength{\parindent}{0pt}
\setlength{\parskip}{5pt}
\setlength{\emergencystretch}{3em}
\newcommand{\N}{\mathbb{N}}
\newcommand{\Z}{\mathbb{Z}}
\newcommand{\Q}{\mathbb{Q}}
\newcommand{\R}{\mathbb{R}}
\newcommand{\C}{\mathbb{C}}
\newcommand{\F}{\mathbb{F}}
\newcommand{\A}{\mathbb{A}}
\newcommand{\PP}{\mathbb{P}}
\newcommand{\E}{\mathbb{E}}
\newcommand{\eps}{\varepsilon}
\newcommand{\dd}{\,\mathrm{d}}
\begin{document}
\section*{UnsolvedMath ID: 5500052}
\textbf{Source problem code:} catalogue entry
\subsection{Definitions and mathematical statement}
Let $G=(V,E)$ be a finite simple undirected graph. A queue layout consists of a linear order $<$ on $V$ and a partition $E=E_1\cup\cdots\cup E_q$ such that no two edges in the same $E_i$ are nested: edges $xy$ and $vw$, labeled with $x<y$ and $v<w$, are nested when $v<x<y<w$ or $x<v<w<y$. Shared endpoints do not constitute strict nesting. The queue-number $\operatorname{qn}(G)$ is the minimum such $q$.

Is there an absolute integer $C$ such that
\[
\operatorname{qn}(G)\leq C\qquad\text{for every finite planar graph }G\,?
\]
Equivalently, can each planar graph's vertices be ordered so that the maximum size of a rainbow, a set of pairwise nested edges, is bounded by a universal constant?

The supplied literature summary reports an affirmative resolution by Dujmović, Joret, Micek, Morin, Ueckerdt, and Wood (2020), with $C=49$. It further reports improvements to $42$ (2021), and to $38$ for general planar graphs and $25$ for bipartite planar graphs in the 2024 work identified as ``From Tripods to Bipods.'' These numerical improvements are retained as supplied historical claims, not new optimality assertions.
\subsection{Short English statement}
Can every planar graph be ordered so that its edges split into a universally bounded number of classes with no nested pair? The supplied 2020 reference answers yes, initially with at most 49 queues.
\subsection{Sources}
\begin{itemize}
\item \emph{The Open Problems Project}, Problem 52, ``Queue-Number of Planar Graphs,'' \url{https://topp.openproblem.net/p52}; accessed 2026-07-29, source-TeX extraction. The background reports its maintained entry as open.
\item Dujmović, Joret, Micek, Morin, Ueckerdt, and Wood, ``Planar Graphs Have Bounded Queue Number,'' \emph{Journal of the ACM} (2020), arXiv:1904.04791, as the affirmative-resolution citation supplied in the record. The summary names the conjecture after Heath--Leighton--Rosenberg.
\item Supplied follow-up attributions: an improvement to 42 (2021), without a full citation, and ``From Tripods to Bipods'' (2024), reported bounds 38 and 25. The record provides no further identifiers for these improvements.
\item TOPP source repository \url{https://github.com/edemaine/topp}; the supplied rights note says its LICENSE expressly scopes software, so problem-text reuse remains NEEDS\_REVIEW.
\item Local registry \texttt{data/unsolvedmath/problems.json}, permanent UnsolvedMath ID $5500052$, source code catalogue entry, \emph{Queue-Number of Planar Graphs}. Dataset attribution: Ulam AI and the UnsolvedMath Contributors.
\end{itemize}
\end{document}
